\SourceLeaf{292}%
Grösse einer Dreiecksseite \(r\), ihrem Azimuthe \(T\) an dem Anfangspunkte, und der Breite dieses Anfangspunkts \(S\), abzuleiten das Azimuth der Seite an dem andern Endpunkte \(T'\pm180^0\), die Breite desselben \(S'\) und den Längenunterschied beider Punkte \(\lambda\). Da dies nichts weiter ist als die Auflösung eines sphärischen Dreiecks, so verdient diese Aufgabe nur deshalb hier einen Platz, weil die gewöhnlich gebrauchten Formeln hier einiger Umformung bedürfen, wenn man in den Resultaten (nach der Bemerkung im 10 Art.) dieselbe Genauigkeit erreichen will, in welcher \(r\) gegeben ist, ohne mehrziffrige Logarithmen zu Hülfe zu nehmen. Um unter den verschiedenen Auflösungsarten nach jedesmaligem Bedürfniss wählen zu können, setze ich zuvörderst diejenigen hieher, die auf den bekannten elementaren Formeln der sphärischen Trigonometrie beruhen.

Erste Methode
\[
\begin{aligned}
\operatorname{tang}s&=\cos T\operatorname{tang}r\\
\operatorname{tang}\lambda&=\frac{\operatorname{tang}T\sin s}{\cos(S-s)}\\
\operatorname{tang}S'&=\cos\lambda\operatorname{tang}(S-s)\\
\sin T'&=\frac{\sin T\cos S}{\cos S'}
\end{aligned}
\]
\indent Zweite Methode
\[
\begin{aligned}
\operatorname{tang}R&=\frac{\operatorname{tang}S}{\cos T}\\
\operatorname{tang}T'&=\frac{\operatorname{tang}T\cos R}{\cos(R-r)}\\
\operatorname{tang}S'&=\cos T'\operatorname{tang}(R-r)\\
\sin\lambda&=\frac{\sin r\sin T}{\cos S'}=\frac{\sin r\sin T'}{\cos S}
\end{aligned}
\]
\indent Dritte Methode
\[
\begin{aligned}
\sin(45^0+\tfrac{1}{2}S')\sin\tfrac{1}{2}(T'+\lambda)&=\sin(45^0+\tfrac{1}{2}(S+r))\sin\tfrac{1}{2}T\\
\sin(45^0+\tfrac{1}{2}S')\cos\tfrac{1}{2}(T'+\lambda)&=\sin(45^0+\tfrac{1}{2}(S-r))\cos\tfrac{1}{2}T\\
\cos(45^0+\tfrac{1}{2}S')\sin\tfrac{1}{2}(T'-\lambda)&=\cos(45^0+\tfrac{1}{2}(S+r))\sin\tfrac{1}{2}T\\
\cos(45^0+\tfrac{1}{2}S')\cos\tfrac{1}{2}(T'-\lambda)&=\cos(45^0+\tfrac{1}{2}(S-r))\cos\tfrac{1}{2}T
\end{aligned}
\]
\indent In Beziehung auf die Kürze der Rechnung hat die dritte Methode einigen Vorzug vor den beiden andern, während diese im Allgemeinen die Resultate ein wenig schärfer geben können, namentlich \(\lambda\) immer mit völlig genügender Schärfe: \(T'\) wird aber, wenn es einem rechten Winkel nahe kommt, durch die erste Methode vergleichungsweise nur ungenau bestimmt. Verlangt man aber alle drei Resultate mit gleichmässiger und, aus dem Gesichtspunkte des 10 Art. betrach%
